% Alexander Tong | NASA Pathways Engineering Resume
% Compile with pdfLaTeX in Overleaf
\documentclass[letterpaper,11pt]{article}

\usepackage[T1]{fontenc}
\usepackage[utf8]{inputenc}
\usepackage{lmodern}
\usepackage[margin=0.5in]{geometry}
\usepackage{titlesec}
\usepackage{enumitem}
\usepackage[hidelinks]{hyperref}
\usepackage[english]{babel}
\input{glyphtounicode}

\pdfgentounicode=1
\pagestyle{empty}
\urlstyle{same}
\raggedright
\raggedbottom
\setlength{\parindent}{0pt}

\titleformat{\section}
  {\large\scshape\raggedright}
  {}{0em}{}[\titlerule]
\titlespacing*{\section}{0pt}{7pt}{4pt}

\newcommand{\resumeEntry}[4]{
  {\small\textbf{#1 \textbar{} #2}}\\[-1pt]
  {\small\textit{#3}}\\[-1pt]
  {\small #4}
}

\newcommand{\resumeItem}[1]{\item \small{#1}}
\newcommand{\resumeItemListStart}{
  \begin{itemize}[
    leftmargin=0.17in,
    label=\textbullet,
    itemsep=1.5pt,
    topsep=2pt,
    parsep=0pt,
    partopsep=0pt
  ]
}
\newcommand{\resumeItemListEnd}{
  \end{itemize}\vspace{3pt}
}

\begin{document}

\begin{center}
  {\Huge\bfseries\scshape Alexander Tong}\\[3pt]
  \small
  860-706-3624
  \enspace\textbar\enspace
  \href{mailto:alto335@bu.edu}{alto335@bu.edu}
  \enspace\textbar\enspace
  \href{https://alexandertong.space}{alexandertong.space}
  \enspace\textbar\enspace
  \href{https://www.linkedin.com/in/alto335/}
  {linkedin.com/in/alto335}
\end{center}
\vspace{-8pt}

\section{Education}

{\small
\textbf{Boston University, College of Engineering}
\textbar{} Boston, MA
\hfill Expected May 2028\\
B.S. Mechanical Engineering, Aerospace Concentration
\textbar{} GPA: 3.71/4.00
}\\
{\small
97 semester credits awarded: 73 from BU coursework and 24 from AP examinations
}

\section{Relevant Engineering Experience}

\resumeEntry
  {Argo Separation and Recovery Subsystem}
  {06/2025 -- Present}
  {Co-Lead Engineer, Boston University Rocket Propulsion Group
   \textbar{} Boston, MA}
  {Part-time student club project
   \textbar{} Total BURPG commitment: 10--30 hours/week, variable}

\resumeItemListStart
  \resumeItem{
    Designed redundant single-separation, dual-deploy recovery for a
    liquid rocket with a 135 lbm dry mass targeting 100,000 ft;
    delivered a final 18.7 lbm system within a 20 lbm budget through
    four iterative design reviews
  }
  \resumeItem{
    Compared system architectures, commercial off-the-shelf (COTS)
    flight computers, and vacuum-tolerant separation mechanisms
    by mass, cost, lead time, and failure modes; finalized system cost
    at \$1,076
  }
  \resumeItem{
    Calculated 823 lbf parachute shock and 1,230 lbf snatch loads
    using five published methods; designed functional
    redundancy with four cable cutters, two in-series tender
    descenders, and two flight computers
  }
  \resumeItem{
    Wrote standard operating procedures (SOPs), test plans, and safety
    protocols for black-powder separation and recovery tests; led operations by assigning crews, briefing personnel,
    authorizing firing and abort decisions, and directing post-test analysis
  }
  \resumeItem{
    Responded to integrated-test failures and parachute damage by
    developing dual-mode black-powder/CO$_2$ deployment hardware;
    performed pressure and Lam\'{e} thick-wall stress calculations and
    contributed CAD, drawings, machining, pressure-data collection,
    and component testing in a three-person team; planned integrated validation
  }
  \resumeItem{
    Built and operated a 352 L vacuum chamber at 0.011 atm using
    salvaged pumps and available stock; acquired and processed pressure
    data using a LabJack, pressure transducer, and two-amplifier signal
    conditioning for deployment-hardware and avionics tests
  }
  \resumeItem{
    Introduced and dynamically tested Brummel eye splices to reduce
    knots and quick-links, confirming a minimum 1.5 factor of safety
    (FOS); redesigned and machined 49.5\%-lighter cable cutters with a
    calculated 2.01 FOS
  }
\resumeItemListEnd

\resumeEntry
  {Boston University Plummer Laboratory}
  {06/2026 -- Present}
  {Mechanical Engineering Intern \textbar{} Boston, MA}
  {06/2026 -- 08/2026: full-time, 40 hours/week
   \textbar{} 09/2026 -- Present: part-time, 10 hours/week}

\resumeItemListStart
  \resumeItem{
    CNC-milled a 6061-T6 shaker adapter on a Haas VF-2; removed
    40.4\% of starting stock mass, from 3.19 to 1.90 kg, using a conical
    underside chamfer, corner fillets, and 12 drilled holes
  }
  \resumeItem{
    Generated G-code, engineering drawings, and tolerance stack-ups
    for $\pm$0.005 in holes; reduced estimated plate-and-fastener cost
    from \$1,170 outsourced to \$71.11 through in-house fabrication
  }
  \resumeItem{
    Predicted 0.44 mm fixture deflection under a 3 kg load at
    20.1 Hz using SolidWorks dynamic finite element analysis (FEA);
    analyzed bolt preload and joint failure modes, including 1.49 and
    2.46 minimum factors of safety
  }
  \resumeItem{
    Designed and executed 48 tests across four physical Ising-model
    designs using an Instron 68SC-5; processed 0--181 mN force data
    and presented methods, hardware, and initial results at a
    multi-laboratory meeting
  }
\resumeItemListEnd

\resumeEntry
  {Culturon}
  {02/2026 -- 06/2026}
  {Mechanical Engineering Intern \textbar{} Sydney, Australia}
  {Part-time \textbar{} 10 hours/week}

\resumeItemListStart
  \resumeItem{
    Planned and executed the company's first full-scale nanoparticle
    deposition map across chamber locations and magnetic configurations;
    used scanning electron microscopy (SEM), optical microscopy, and
    ImageJ to characterize polarity-dependent deposition and asymmetric flow
  }
  \resumeItem{
    Formed a working hypothesis for contamination cause and pattern
    from sample-plate SEM/ImageJ results, optical microscopy, and direct
    visual evidence
  }
  \resumeItem{
    Presented pump-port/chamber and electron-curtain prevention concepts
    based on the hypothesis and fluid/electromagnetic first principles
    in a technical design review
  }
  \resumeItem{
    Designed Onshape sheet-metal sample racks and coupons for
    reactor-port insertion and adhesive-free mounting with bend reliefs,
    fabrication tolerances, assembly clearances, and limited flow disturbance
  }
\resumeItemListEnd

\newpage

\section{Selected Engineering Projects}

\resumeEntry
  {Aerodynamic Downcomer Fairing}
  {08/2026 -- Present}
  {Project Lead, Boston University Rocket Propulsion Group
   \textbar{} Boston, MA}
  {Part-time student club project
   \textbar{} Included in total BURPG commitment listed above}

\resumeItemListStart
  \resumeItem{
    Leading set-based design of an aerodynamic and composite fairing to protect
    Argo's external downcomer and wiring; evaluating leading-edge,
    midsection, and boattail profiles for Mach 2.81 max-Q flow using
    wave-, skin-friction-, and base-drag considerations
  }
  \resumeItem{
    Developing Ansys computational fluid dynamics (CFD) models at Mach
    2.0, 2.5, and 2.81 and applying adjoint optimization under packaging
    and manufacturability constraints; directing an independent
    SolidWorks Flow Simulation comparison
  }
\resumeItemListEnd

\resumeEntry
  {AIAA Design-Build-Fly Aircraft}
  {06/2025 -- 12/2025}
  {Wing Structural Lead, Boston University \textbar{} Boston, MA}
  {Part-time student club project
   \textbar{} 10--20 hours/week, variable}

\resumeItemListStart
  \resumeItem{
    Led two wing designs in Fusion 360; produced an
    extruded-polystyrene (XPS) prototype 50\% lighter than the prior
    year's wing and cut build time from three months to one
  }
  \resumeItem{
    Designed a split NACA 2412 wing with 0.52 taper ratio and
    5-degree dihedral from aerodynamic-team specifications and
    first-prototype manufacturing lessons; developed the servo flap
    system
  }
  \resumeItem{
    Led structural and manufacturing reviews, mentored nine members
    in laser cutting, waterjet operation, and Fusion 360, and trained
    the succeeding wing lead
  }
\resumeItemListEnd

\resumeEntry
  {Hybrid Rocket Nozzle}
  {09/2024 -- 11/2024}
  {Lead Engineer, Boston University Rocket Propulsion Group
   \textbar{} Boston, MA}
  {Part-time student club project
   \textbar{} Approximately 5 hours/week}

\resumeItemListStart
  \resumeItem{
    Led five members through three design reviews to develop a
    7.86 lbf nozzle optimized for 45,000 ft using isentropic flow
    relations, an 80\%-length Rao contour, and SolidWorks CAD/FEA
  }
  \resumeItem{
    Reduced mass 73\% from the initial design, increased predicted
    thrust 5\%, and calculated a 7.8 FOS under mechanical loads
    alone; incorporated review feedback on manufacturability and fragility
  }
  \resumeItem{
    Led the team that designed, soldered, and programmed the launch
    computer and completed an approximately two-second hotfire;
    identified a load-cell malfunction that prevented thrust-data
    collection
  }
\resumeItemListEnd

\resumeEntry
  {Rocket Test Stand Fire Suppression System}
  {02/2025 -- 05/2025}
  {Project Contributor, Boston University Rocket Propulsion Group
   \textbar{} Boston, MA}
  {Part-time student club project
   \textbar{} Included in total BURPG commitment listed above}

\resumeItemListStart
  \resumeItem{
    Planned fabrication and designed the SolidWorks assembly,
    mounting, adapter pieces, 3D-printed mounts, and sheet-metal
    fixtures for a remote CO$_2$ fire-suppression system targeting
    0.23 kg/s flow
  }
  \resumeItem{
    Developed stainless-steel pipe routing under SAE AS130B guidance
    and translated project-supplied requirements of 1,070 psi and a
    minimum 4.0 FOS into the mechanical layout
  }
\resumeItemListEnd

\resumeEntry
  {Ablative Anode Pulsed Plasma Thruster Laboratory}
  {05/2026 -- 06/2026}
  {Volunteer, University of Sydney \textbar{} Sydney, Australia}
  {Part-time volunteer work
   \textbar{} Approximately 5 hours/week}

\resumeItemListStart
  \resumeItem{
    Designed an acrylonitrile butadiene styrene (ABS) inductor spool,
    performed coil-winding calculations, checked vacuum-outgassing
    suitability, and supported the established thruster operating
    procedure
  }
\resumeItemListEnd

\section{Technical Skills and Recognition}

{\small
\textbf{Design and analysis:}
SolidWorks, Fusion 360, Onshape, Ansys CFD, adjoint optimization,
compressible-flow analysis, FEA, geometric dimensioning and
tolerancing (GD\&T), engineering drawings, MATLAB, Python\\
\textbf{Fabrication and testing:}
CNC and manual machining, composite layup, 3D printing, vacuum
systems, LabJack data acquisition, pressure instrumentation, test
planning and operations\\
\textbf{Recognition:}
Boston University Distinguished Summer Research Fellowship, 2024
}

\end{document}
